\chapter{Reform, and its reversal}
\label{ch:dollarreform}

\section{Reform under duress}

Between July 1993 and the end of 1995 the Cuban government did more to
liberalise its economy than in the previous thirty-four years combined. It did
so reluctantly, described the measures throughout as temporary concessions
forced by circumstance, and withdrew as many of them as it could as soon as it
was able. That sequence is the subject of this chapter, and it is the clearest
single instance of the pattern this book has been assembling.

The measures, in order:

\emph{July 1993.} Possession of hard currency was decriminalised. Until that
month a Cuban caught with a dollar could be prosecuted; from that month Cubans
could hold, spend and receive dollars legally. This was the pivotal decision,
because it legalised remittances --- and remittances from the diaspora that
thirty years of policy had denounced as worms became, within a decade, one of
the largest sources of foreign exchange in the country.

\emph{September 1993.} Self-employment was re-legalised in 117 categories, later
extended to around 157. The list was a list: activity was permitted by
enumeration, not by right, and university-trained professions were excluded ---
a doctor could not practise privately, an engineer could not consult, but either
could rent a room or drive a taxi. This design decision, more than any other,
determined what the Cuban private sector became.

\emph{September 1993.} The state farms, which held most of the agricultural
land, were broken into Basic Units of Cooperative Production, the UBPCs, holding
land in usufruct with a nominal right to organise their own work. Roughly 40 per
cent of agricultural land changed hands on paper.\unv{} In practice the state
retained the procurement monopoly, the input supply and the price, so what was
devolved was the loss and not the decision. The UBPCs have underperformed the
private and older-cooperative sector continuously ever since.

\emph{October 1994.} Agricultural markets were reopened at free prices for
produce above the compulsory state quota --- the same measure that had been
introduced in 1980 and abolished in 1986. Food appeared in Havana within weeks,
at prices most people could not afford, which was also what had happened in
1980.

\emph{1994--95.} The fiscal and monetary stabilisation, which is the part of
this period that gets least attention and deserves the most. The budget deficit,
which had reached something like 30 per cent of output in 1993, was cut to
around 7 per cent in 1994 and to 2 per cent by 1996, through price rises,
subsidy cuts, new taxes under a 1994 tax law, and the removal of currency by
confiscatory means. Exchange houses were opened in 1995 at a market rate. The
informal exchange rate, which had reached 120 to 150 pesos to the dollar in mid-
1994, came back to around 20 to 25 by 1996 and stayed in that band for the next
twenty-five years \citep{mesalago2005}.\unv{} A currency that had been destroyed was
stabilised inside two years without external assistance. Chapter~\ref{ch:money} takes this
seriously as a precedent, because the same state did the opposite in 2021.

\emph{September 1995.} Law 77 opened foreign investment, permitting up to full
foreign ownership in principle, and creating the joint-venture form that most
foreign capital in Cuba has used since. It kept the feature that has done most
to deter investors: the foreign partner does not hire its own workers, but
contracts them from a state employment agency, pays the agency in hard currency,
and the agency pays the worker in pesos, keeping the difference. At an official
rate of one to one, this meant a worker whose employer paid the state several
hundred dollars a month received the peso equivalent of perhaps twenty.
Chapter~\ref{ch:capital} argues that this single arrangement has cost Cuba more
foreign investment than any American statute.

The measures worked. 1994 was the trough; growth resumed in 1995 and ran through
the second half of the decade. The economy that emerged, however, was not the
one that had gone in.

\section{The dual economy}

Cuba came out of the Special Period with two currencies, two price levels, two
labour markets and two populations.

The convertible peso, the CUC, was introduced in 1994 and pegged at one to the
dollar. Alongside it the ordinary peso continued, at twenty-four or twenty-five
to the CUC for households and --- this is the part that mattered --- at one to
one for state enterprise accounting. Cuba thus ran a twenty-five-fold
discrepancy between the rate its households faced and the rate its own firms
used to value imports, exports, costs and profits, for twenty-six years. Every
enterprise account in the country was fiction. No Cuban manager could tell
whether their operation created or destroyed value, and neither could the
ministry above them. This is not a technicality; it is the reason Cuban
industry could not be reformed even by people who wanted to reform it, because
nobody could identify which parts were worth keeping.

The social consequence was the inversion of the pyramid. A surgeon earned a
state salary in pesos. A bellhop earned tips in convertible currency. A
professor of literature drove a taxi in the afternoon; a research chemist rented
her spare room. The valuation the society had spent thirty years teaching ---
that education was the route to standing --- was contradicted daily and
publicly by the exchange rate. Cuba's Gini coefficient, which had been around
0.24 in the mid-1980s and was among the lowest measured anywhere, was around
0.40 by the end of the 1990s \citep{mesalago2005}.\unv{}

That inequality had a colour. The emigration of 1959--62 had been
overwhelmingly white; the remittances that flowed back thirty years later
therefore flowed overwhelmingly to white households, by margins that the
available surveys put at four fifths and above.\unv{} The tourist sector, the
other route to hard currency, hired front-of-house staff by \emph{buena
presencia} --- good appearance --- which in practice meant what it has always
meant. So the state that had abolished the colour bar in 1959 and then declared
the racial question closed in 1962, and had thereby made it unmeasurable for
thirty years, entered its market opening with foreign-currency access
distributed along the exact line it had stopped looking at.
Chapter~\ref{ch:premises} treats this as a hard constraint on any transition
design, and not as a historical note.

Tourism was the visible face of it. Arrivals went from about 340,000 in 1990 to
around 745,000 in 1995 and to about 1.8 million by 2000.\unv{} The plant was
built by joint ventures with Spanish, Canadian and Jamaican chains, mostly as
all-inclusive beach resorts at Varadero and the northern cays, physically
separated from the country around them --- and, until 2008, legally separated
too, since Cubans were not permitted in tourist hotels. A state that had made
its founding claim on the abolition of a foreigners' Cuba had rebuilt one, and
this time it was the government enforcing the door.

\section{The reversal}

By 1996 the emergency had passed, and the leadership's assessment of the
reforms changed with it.

The signal came in March 1996, when Ra\'ul Castro delivered a report to the
Central Committee attacking the Centre for the Study of the Americas, a party
research institute whose economists had been analysing the reforms and their
possible extension. The institute was dismantled and its researchers dispersed.
The message to Cuban social science was received and has never needed
repeating.

What followed was a decade of quiet reversal. Self-employment licences, which
had peaked at around 209,000 in early 1996, were squeezed by taxation, by
inspection, and by the withdrawal of categories, falling to roughly 150,000 by
2003 and lower after \citep{ritter2015}.\unv{} New licences in most categories were suspended
outright in 2003--04. In 2003 enterprises were required to hold and transact
their foreign currency through a single central account, ending the
decentralised access that had made the joint-venture sector function. In
November 2004 the United States dollar was withdrawn from circulation
altogether and a 10 per cent surcharge imposed on converting dollars into CUC,
which was presented as a response to American financial pressure and which also
happened to tax every remittance.

None of this was a return to 1968; the markets stayed open, the tourism sector
grew, the remittances kept coming. But the direction was unmistakable, and its
logic was consistent: the reforms of 1993--95 had been concessions extracted by
an emergency, and with the emergency over the concessions were withdrawn.

\section{Politics in the same years}

The reversal was not only economic, and the external environment did not help.

On 24 February 1996 Cuban air-force MiGs shot down two civilian aircraft flown
by Brothers to the Rescue, an exile organisation that had been searching for
rafters and dropping leaflets over Havana, killing four men. The location of the
shootdown --- inside or outside Cuban airspace --- was disputed; that the
aircraft were unarmed civilian light planes was not. The political consequence
in Washington was immediate: the Cuban Liberty and Democratic Solidarity Act,
Helms--Burton, which had been stalled, passed within weeks. It codified the
embargo into statute, so that it can now be lifted only by Congress and not by a
president; it extended sanctions to foreign companies trafficking in confiscated
American property; and its Title III created a right for American nationals to
sue such companies in American courts. Every president from Clinton onward
suspended Title III at six-month intervals until 2019.

Helms--Burton matters to Part III for a reason beyond its direct effect. It
converted American policy toward Cuba from an executive instrument, which can be
traded, into a legislative one, which must be repealed. Any Cuban reform design
that assumes the embargo will end when relations improve is assuming something
the American constitutional structure does not readily deliver, and
Chapter~\ref{ch:premises} therefore treats the embargo as a variable that may
not move.

Inside Cuba the same years produced Law 88 of 1999, which made collaboration
with United States policy punishable by up to twenty years and which is
sufficiently broadly drafted to cover most forms of independent journalism; and
then, in 2002, the Varela Project. Oswaldo Pay\'a and a network of Catholic lay
activists used Article 88 of the 1976 constitution, which permitted citizens to
propose legislation with 10,000 signatures, to submit 11,020 signatures asking
for a referendum on free expression, free association, an amnesty and free
elections. It was the most serious attempt to use the system's own rules against
it in forty years. The response was a counter-mobilisation that collected some
eight million signatures for a constitutional amendment declaring the socialist
character of the system irrevocable, which the National Assembly duly passed in
June 2002. The route Pay\'a had found was closed by constitutional amendment.

In March 2003, in the week the United States invaded Iraq, Cuban state security
arrested seventy-five independent journalists, librarians, trade unionists and
opposition organisers, tried them within weeks, and sentenced them to terms of
six to twenty-eight years. It is known as the Black Spring. In the same month
three young men hijacked a Havana Bay passenger ferry in an attempt to reach
Florida, took hostages, harmed nobody, surrendered when the boat ran out of
fuel, and were tried and executed by firing squad within nine days. Cuba's
standing with the European left, which had survived Padilla, largely did not
survive this. The wives and mothers of the seventy-five began walking to mass in
white every Sunday, and were still doing so, and being detained for it, twenty
years later.

\begin{ledger}{reform and reversal, 1993--2004}
\emph{Achieved.} An economy pulled out of free fall in eighteen months. A
destroyed currency stabilised without external assistance --- deficit from
about 30 per cent of output to 2 per cent, informal exchange rate from 150 to
25 --- which is a genuine technocratic accomplishment and the counter-example
to Chapter~\ref{ch:crisis}. Dollars, self-employment, farmers' markets and
foreign investment legalised. Tourism rebuilt from nothing to 1.8 million
visitors. A private sector of a few hundred thousand people re-created from
zero, and with it the first Cubans in a generation who did not depend on the
state for their income.

\emph{Cost.} A dual currency that made every enterprise account in the country
fictitious for twenty-six years, and thereby made the economy unreformable by
anyone, including those who wanted to reform it. Inequality up from among the
lowest measured anywhere to Latin American levels in a decade, distributed
along a racial line the state had spent thirty years declining to measure.
Tourist enclaves Cubans were legally barred from entering. A reform programme
withdrawn as soon as it was survivable to withdraw it, with the research
institute that had analysed it dismantled as a warning. Seventy-five people
imprisoned for terms up to twenty-eight years, and three shot, in a single
month.
\end{ledger}
